\documentclass[aps,pre,reprint,amsmath,amssymb,superscriptaddress,nofootinbib]{revtex4-2}

\usepackage{bm}
\usepackage{booktabs}
\usepackage{graphicx}

\newcommand{\RR}{\mathbb{R}}
\newcommand{\dd}{\mathrm{d}}
\newcommand{\E}{\mathbb{E}}
\newcommand{\Var}{\operatorname{Var}}
\newcommand{\Cov}{\operatorname{Cov}}
\newcommand{\grad}{\nabla}
\newcommand{\Tr}{\operatorname{Tr}}
\newcommand{\cA}{\mathcal{A}}
\newcommand{\cH}{\mathcal{H}}
\newcommand{\cW}{\mathcal{W}}
\newcommand{\cG}{\mathcal{G}}
\newtheorem{theorem}{Theorem}

\begin{document}

\title{Matched-equilibrium lower bound for polynomial current noise in a torque-driven underdamped Gaussian particle}

\author{Xiaojian Zeng}
\email{xiaojianzeng@foxmail.com}
\affiliation{PBC School of Finance, Tsinghua University, 100190, China}

\date{September 12, 2026}

\begin{abstract}
An equilibrium process with the same position marginal as a nonequilibrium
underdamped diffusion provides a natural reference for current fluctuations,
but the reference need not order arbitrary observables.  We resolve this
comparison for a stable two-dimensional Ornstein--Uhlenbeck particle driven by
a constant torque and for the complete family of position-current weights
whose homogeneous degrees are one and three.  Gradient weights are stationary
boundary increments, leaving a four-dimensional quotient of the original
twelve-dimensional coefficient space.  After an exact nondimensionalization,
rotational covariance splits that quotient into a coupled isotropic block and
an angular-two doublet.  Complete polynomial Poisson equations of total degree
at most four give the driven-minus-equilibrium diffusion matrix as rational
functions of two positive parameters.  Its isotropic leading minor and
determinant and its angular coefficient have explicit strictly positive
certificates over the full stability domain.  Consequently, every current in
the stated family is at least as noisy in the driven process as in the matched
equilibrium trap, with equality exactly for gradient boundary currents.  The
calculation is exact for the linear dynamics and is independently checked by a
two-time covariance and Wick-contraction representation.
\end{abstract}

\maketitle

\section{Introduction}
\label{sec:introduction}

Current fluctuations quantify the precision of transport in nonequilibrium
steady states.  In underdamped systems their behavior is especially sensitive
to the observable and to forces that act on velocity or break time-reversal
symmetry.  For example, a magnetic field can invalidate familiar
thermodynamic-uncertainty bounds for work and heat currents
\cite{ChunFischerSeifert2019}.  It is therefore useful to separate statements
about a particular current from statements that order a whole family of
observables.

Dechant introduced an ``equivalent equilibrium'' dynamics for an underdamped
steady state: the reference retains the position marginal while removing the
local mean flow.  In a torque-driven Brownian-gyrator sector, the torque-work
current is noisier than its equivalent-equilibrium counterpart, and the same
ordering was conjectured for more general observables
\cite{Dechant2022}.  The conjecture cannot hold without specifying both the
model and the observable class.  Recent velocity-resolved examples show that
closely related fluctuation inequalities can fail when the readout or physical
premises change \cite{VuVoBaoSaito2026}.

Here we give an all-parameter result for a finite but complete polynomial
sector of the zero-field torque-driven Gaussian model.  The weight of the
position current may contain every real homogeneous term of degree one and
three.  This sector is already beyond the linear current-type observables in
the exact large-deviation treatment of linear diffusions by du Buisson and
Touchette \cite{duBuissonTouchette2023}; cubic position weights produce
quartic phase-space rates.  Their erratum concerns eigenfunction
normalizability and full scaled-cumulant-generating-function domains
\cite{duBuissonTouchette2025Erratum}, whereas our argument uses only the
centered Poisson equation at zero tilt.  Results for occupation estimators in
nonreversible samplers also have different matching premises: those processes
share the full invariant law, not only the position marginal
\cite{CoghiChetriteTouchette2021}.

Our result has three ingredients.  First, quotienting gradient boundary
increments reduces the twelve current coefficients to four without discarding
any extensive current.  Second, a change of units leaves only two positive
dimensionless parameters and preserves the sign of the diffusion difference
by a nonsingular congruence.  Third, rotational covariance reduces the exact
degree-four Poisson calculation to a two-dimensional invariant block and an
angular-two doublet.  The resulting global sign certificate proves the
ordering throughout the strict stability domain, rather than at sampled
parameters.  The closest sources identified above do not state this cubic
quotient theorem; that comparison is necessarily a bounded priority search,
not a claim of exhaustive bibliographic precedence.

\section{Model, stationary laws, and current space}
\label{sec:model}

Let
\begin{equation}
 R=\begin{pmatrix}0&-1\\1&0\end{pmatrix},
 \qquad x,v\in\RR^2.
\end{equation}
The driven process is
\begin{subequations}
\label{eq:driven}
\begin{align}
 \dd x_t &=v_t\,\dd t,\\
 m\,\dd v_t&=(-k x_t+\kappa R x_t-\gamma v_t)\,\dd t
       +\sqrt{2\gamma T}\,\dd W_t,
\end{align}
\end{subequations}
where $W_t$ is standard two-dimensional Brownian motion.  We assume
\begin{equation}
 m,\gamma,k,T>0,\qquad \kappa\ne0,\qquad
 k_*:=k-\frac{m\kappa^2}{\gamma^2}>0.
\label{eq:domain}
\end{equation}
The comparison process keeps $m,\gamma,T$ but replaces the entire position
force with the conservative force $-k_*x$:
\begin{subequations}
\label{eq:equilibrium}
\begin{align}
 \dd x_t &=v_t\,\dd t,\\
 m\,\dd v_t&=(-k_* x_t-\gamma v_t)\,\dd t
       +\sqrt{2\gamma T}\,\dd W_t.
\end{align}
\end{subequations}
The matched stiffness is $k_*$, not $k$.

Both processes are initialized in their own complete stationary laws.  Direct
solution of the Lyapunov equation gives, for Eq.~\eqref{eq:driven},
\begin{equation}
 x\sim N\!\left(0,\frac{T}{k_*}I_2\right),\qquad
 v\mid x\sim N\!\left(\frac{\kappa}{\gamma}Rx,
                       \frac{T}{m}I_2\right),
\label{eq:stationary-driven}
\end{equation}
whereas Eq.~\eqref{eq:equilibrium} has
\begin{equation}
 x\sim N\!\left(0,\frac{T}{k_*}I_2\right),\qquad
 v\sim N\!\left(0,\frac{T}{m}I_2\right),\qquad x\perp v.
\label{eq:stationary-equilibrium}
\end{equation}
Thus the position marginals agree although the phase-space laws do not.  For
each complex spatial mode the deterministic characteristic equation is
$m\lambda^2+\gamma\lambda+k\mp i\kappa=0$; its strict Hurwitz condition is
$m\kappa^2<\gamma^2k$, exactly Eq.~\eqref{eq:domain}.  There is no spatial
cutoff or boundary other than Gaussian decay at infinity.  We set Boltzmann's
constant to one, so $T$ has units of energy, while $k$ and $\kappa$ have
units of stiffness.

Let $\cH_d$ denote the real homogeneous scalar polynomials of degree $d$ in
$(x_1,x_2)$, and define the twelve-dimensional current-weight space
\begin{equation}
 \cW_{1,3}=(\cH_1\oplus\cH_3)^2.
\end{equation}
For $w\in\cW_{1,3}$, use the same weight in both processes and set
\begin{equation}
 J_w(t)=\int_0^t w(x_s)\mathbin{\cdot}v_s\,\dd s,
 \qquad
 D(w)=\lim_{t\to\infty}\frac{\Var J_w(t)}{2t}.
\label{eq:current-diffusion}
\end{equation}
The two diffusion rates will be denoted $D_{\rm d}$ and $D_{\rm e}$.
Because $\dd x=v\,\dd t$ has finite variation, Eq.~\eqref{eq:current-diffusion}
has no It\^o--Stratonovich ambiguity.  The variance is centered by definition;
the two current means need not agree.

\section{Gradient quotient and statement of the result}
\label{sec:quotient}

For any $\phi\in\cH_2\oplus\cH_4$,
\begin{equation}
 J_{\grad\phi}(t)=\phi(x_t)-\phi(x_0).
\end{equation}
Stationarity and Gaussian moments make its variance bounded uniformly in
$t$.  Its covariance with an extensive current is $o(t)$ by
Cauchy--Schwarz.  Hence both diffusion forms are invariant under
$w\mapsto w+\grad\phi$, and the relevant space is
\begin{equation}
 \cW_{1,3}/\cG,\qquad
 \cG=\grad(\cH_2\oplus\cH_4).
\end{equation}
The gradient space has dimension $3+5=8$, leaving a four-dimensional
quotient.

Every quotient class has the unique representative
\begin{equation}
 w_c(x)=\left(c_0+c_1x_1^2+c_2x_1x_2+c_3x_2^2\right)Rx.
\label{eq:quotient-representative}
\end{equation}
Indeed, in two dimensions curl has the gradient fields as its kernel, while
\begin{equation}
 \operatorname{curl}(pRx)=2p+x\mathbin{\cdot}\grad p.
\label{eq:curl}
\end{equation}
It acts as multiplication by $2$ on constant $p$ and by $4$ on
quadratic homogeneous $p$.  Dimension counting then proves existence and
uniqueness in Eq.~\eqref{eq:quotient-representative}; further details are in
Appendix~\ref{app:quotient}.

\begin{theorem}
\label{thm:main}
For every parameter tuple satisfying Eq.~\eqref{eq:domain} and every real
$w\in\cW_{1,3}$,
\begin{equation}
 D_{\rm d}(w)\ge D_{\rm e}(w).
\label{eq:main-inequality}
\end{equation}
Equality holds if and only if
$w\in\grad(\cH_2\oplus\cH_4)$.
\end{theorem}

The theorem concerns only the stationary long-time second-cumulant rate.  It
does not assert a finite-time inequality, a full generating-function ordering,
or an ordering for velocity-dependent or higher-degree weights.

\section{Dimensionless exact Poisson representation}
\label{sec:poisson}

Introduce dimensionless time and length
\begin{equation}
 \theta=\frac{\gamma t}{m},\qquad
 \ell=\frac{\sqrt{mT}}{\gamma},\qquad
 X=\frac{x}{\ell},\qquad V=\frac{mv}{\gamma\ell},
\end{equation}
and the two parameters
\begin{equation}
 a=\frac{mk_*}{\gamma^2}>0,\qquad
 b=\frac{m\kappa}{\gamma^2}\ne0,\qquad u=b^2>0.
\label{eq:dimensionless-parameters}
\end{equation}
The backward generators become
\begin{subequations}
\label{eq:generators}
\begin{align}
 L_b={}&V\mathbin{\cdot}\grad_X+
 [-(a+b^2)X+bRX-V]\mathbin{\cdot}\grad_V+\Delta_V,\\
 L_0={}&V\mathbin{\cdot}\grad_X+[-aX-V]\mathbin{\cdot}\grad_V+\Delta_V.
\end{align}
\end{subequations}
Their covariance blocks are
\begin{equation}
 \begin{aligned}
 \langle XX^{\mathsf T}\rangle&=a^{-1}I,\\
 \langle VX^{\mathsf T}\rangle&=(b/a)R,\\
 \langle VV^{\mathsf T}\rangle&=((a+b^2)/a)I
 \end{aligned}
\label{eq:cov-blocks}
\end{equation}
for the driven process and $(a^{-1}I,0,I)$ for equilibrium.

Let $S_\ell=\operatorname{diag}(1,\ell^2,\ell^2,\ell^2)$ act on the
coefficient vector in Eq.~\eqref{eq:quotient-representative}.  Direct
substitution in the integrated current gives
\begin{equation}
 D_{\rm physical}(c)=\frac{\gamma\ell^4}{m}
 D_{\rm dimensionless}(S_\ell c).
\label{eq:physical-scaling}
\end{equation}
The prefactor is positive and $S_\ell$ is nonsingular.  Thus sign and kernel
dimension are preserved.  Torque reversal is implemented by a spatial
reflection; the final symmetric form depends on $b$ only through $u=b^2$.

For a centered polynomial rate
$\widehat f=f-\langle f\rangle$, solve
\begin{equation}
 -Lh=\widehat f.
\label{eq:poisson}
\end{equation}
Exponential stability makes the centered solution unique up to a constant,
and the Green--Kubo formula gives
\begin{equation}
 D(f)=\langle \widehat f h\rangle.
\label{eq:green-kubo}
\end{equation}
For several basis currents the quadratic-form matrix is the symmetric part of
$\langle\widehat f_i h_j\rangle$.  Because $L_b$ and $L_0$ map
phase-space polynomials of total degree at most four into the same space, the
calculation is an exact finite closure, not a truncation of the dynamics.

Writing
\begin{equation}
 \cA=X_1V_2-X_2V_1,
\end{equation}
the rates associated with Eq.~\eqref{eq:quotient-representative} are a
quadratic prefactor times $\cA$.  Rotational covariance separates the
invariant span of $\cA$ and $|X|^2\cA$ from a real angular-two doublet.
The complete invariant Poisson sector contains 13 nonconstant monomials and
the complex angular-two sector contains 11.  Appendix~\ref{app:closure} lists
these sectors and their stationary means.

\section{Exact quotient matrix and positivity certificate}
\label{sec:certificate}

Decompose the quadratic prefactor as
\begin{equation}
 p=c_0+s|X|^2+d(X_1^2-X_2^2)+eX_1X_2,
\label{eq:p-decomposition}
\end{equation}
where
\begin{equation}
 s=\frac{c_1+c_3}{2},\qquad
 d=\frac{c_1-c_3}{2},\qquad e=c_2.
\end{equation}
The exact driven-minus-equilibrium calculation gives
\begin{equation}
 \Delta D=A c_0^2+2B c_0s+C s^2+Gd^2+\frac{G}{4}e^2.
\label{eq:block-form}
\end{equation}
Define
\begin{subequations}
\label{eq:isotropic-components}
\begin{align}
 A&=\frac{2u(5a+4u+1)}{a^3},\\
 B&=\frac{16u(4a+4u+1)}{a^4},\\
 H&=16a^2+32au+24a+16u^2+40u+9,\\
 \Delta_0&=\frac{32u^2P_0(a,u)}{a^8(4a+3)H},\\
 C&=\frac{\Delta_0+B^2}{A}.
\end{align}
\end{subequations}
The angular coefficient is
\begin{equation}
 G=\frac{16uP_2(a,u)}{
 a^3(4a+3)(4a+3u+3)((a+u)^2+u)K},
\label{eq:angular-coefficient}
\end{equation}
with
\begin{equation}
 K=(16a^2+32au+12a+16u^2-13u)^2
   +4u(20a+20u+3)^2.
\label{eq:K}
\end{equation}
The integer polynomials $P_0$ and $P_2$ are displayed in
Appendix~\ref{app:polynomials}.  All 20 coefficients of $P_0$ and all 33
coefficients of $P_2$ are strictly positive.

In the ordered quotient basis
$(1,X_1^2,X_1X_2,X_2^2)RX$, Eq.~\eqref{eq:block-form} is the matrix
\begin{equation}
 Q_\Delta=
 \begin{pmatrix}
 A&B/2&0&B/2\\
 B/2&(C+G)/4&0&(C-G)/4\\
 0&0&G/4&0\\
 B/2&(C-G)/4&0&(C+G)/4
 \end{pmatrix}.
\label{eq:quotient-matrix}
\end{equation}
Every factor in the denominators of Eqs.~\eqref{eq:isotropic-components} and
\eqref{eq:angular-coefficient} is positive when $a,u>0$.  In particular,
$K>0$ follows immediately from Eq.~\eqref{eq:K}.  Positive polynomial
coefficients imply $P_0,P_2>0$.  Consequently
\begin{equation}
 A>0,\qquad
 \det\begin{pmatrix}A&B\\B&C\end{pmatrix}=\Delta_0>0,
 \qquad G>0.
\label{eq:sylvester}
\end{equation}
Sylvester's criterion makes the isotropic block positive definite, while the
angular-two doublet is strictly positive.  Hence $Q_\Delta$ is positive
definite on the four-dimensional quotient for every $a,u>0$.  The scaling
in Eq.~\eqref{eq:physical-scaling} transfers this statement to every physical
tuple in Eq.~\eqref{eq:domain}.  Gradient classes have zero long-time rate in
both processes, and every nongradient class has $\Delta D>0$.  This proves
Theorem~\ref{thm:main}, including its equality statement.

\begin{figure*}[t]
 \centering
 \includegraphics[width=0.92\textwidth]{figures/quotient_spectrum.pdf}
 \caption{Numerical visualization of the exact quotient formula on an
 $81\times81$ logarithmic grid with $10^{-2}\le a,u\le10^2$.
 Left: smallest eigenvalue of the thermally normalized
 $\widetilde Q_\Delta$.  Right: its spectral condition number.  Every sampled
 eigenvalue is positive, but the grid is illustrative only; positivity on the
 continuous domain follows from Eqs.~\eqref{eq:isotropic-components}--
 \eqref{eq:sylvester}.  The plotted values and regeneration script accompany
 the manuscript.  OpenAI Codex assisted in writing and checking the
 deterministic plotting code; the plotted data were generated from the exact
 formulas, as detailed in Sec.~\ref{sec:reproducibility}.}
 \label{fig:spectrum}
\end{figure*}

\section{Benchmark and independent checks}
\label{sec:checks}

Before applying the closure to cubic currents, we reproduced the published
zero-field torque-work benchmark.  With
$W(t)=\kappa\int_0^t(Rx_s)\cdot v_s\,\dd s$, the source convention gives
\begin{subequations}
\label{eq:work-benchmark}
\begin{align}
 \langle\dot W\rangle&=\frac{2\kappa^2T}{\gamma k_*},\\
 D_W&=\frac{2\kappa^2T^2}{\gamma k_*^3}
 \left(k^2+\kappa^2+\frac{3mk\kappa^2}{\gamma^2}\right),
\end{align}
\end{subequations}
in agreement with the zero-magnetic-field specialization of
Ref.~\cite{ChunFischerSeifert2019}.  That reference uses the opposite torque
orientation, $\kappa_{\rm CFS}=-\kappa$; Eq.~\eqref{eq:work-benchmark} is even
in torque.  Independent real-coordinate Poisson and complex-coordinate
two-time covariance reconstructions reproduce both formulas.

The global proof is symbolic.  Four Poisson identities---two invariant and two
angular, driven and equilibrium---have zero polynomial residual.  The
positive-coefficient audit gives minimum integer coefficients $27$ for
$P_0$ and $1620$ for $P_2$.  Complete real polynomial solves with 69
nonconstant unknowns agree exactly with Eq.~\eqref{eq:quotient-matrix} at the
rational points
\begin{equation}
 (a,b)=(1,1),\quad(2,1/2),\quad(3/2,-2/3),
\end{equation}
including all constant Poisson equations.  These hold-outs test the reduction
and torque sign but do not replace the continuous-domain certificate.

As a separate numerical representation, we form the full two-time covariance
of $(X,V)$, evaluate degree-eight Gaussian moments by Wick contractions, and
integrate symmetrized correlations on $[0,\infty)$.  Across the three physical
cases listed in the reproducibility data, the largest Lyapunov residual is
$5.77\times10^{-14}$, the largest reported quadrature error is
$5.99\times10^{-9}$, and the largest Poisson--Wick relative matrix
discrepancy is $1.053\times10^{-12}$.  These numerical values are auxiliary
error controls; the theorem itself rests on exact identities.

Figure~\ref{fig:spectrum} illustrates the thermally normalized quotient
spectrum.  To put linear and cubic coefficients on comparable stationary
scales, write $c=(\widetilde c_0,a\widetilde c_1,
a\widetilde c_2,a\widetilde c_3)$, so
$\widetilde Q_\Delta=\operatorname{diag}(1,a,a,a)Q_\Delta
\operatorname{diag}(1,a,a,a)$.  The wide variation of its condition number
also explains why a floating-point scan alone would be a poor global sign
test.

\section{Physical interpretation and limitations}
\label{sec:discussion}

The equilibrium reference in Eq.~\eqref{eq:equilibrium} is matched only at the
level of the position marginal.  The driven phase-space law retains a
rotational conditional velocity, Eq.~\eqref{eq:stationary-driven}.  Theorem
\ref{thm:main} states that, despite this distinction, eliminating the rotation
cannot increase the long-time noise of any degree-one/three position current.
The result includes anisotropic cubic readouts: none can undercut the matched
equilibrium calibration.  Equality has a concrete kinematic meaning, namely a
nonextensive endpoint difference $\phi(x_t)-\phi(x_0)$, rather than a hidden
finite-current saturation.

The conclusion is deliberately local in model and observable space.  It uses
linear Gaussian dynamics, zero magnetic field, stationary initialization,
strict stability, and the second cumulant at long times.  It does not establish
the general-observable conjecture of Ref.~\cite{Dechant2022}.  It also does not
apply at the stability boundary, to nonlinear forces, to velocity-dependent
weights, to finite-time variances, to higher cumulants, or to full
large-deviation domains.  In particular, the integrability corrections in
Ref.~\cite{duBuissonTouchette2025Erratum} caution against promoting a local
Poisson calculation to an unproved statement about the complete scaled
cumulant generating function.

The quotient construction and rotational-sector decomposition are reusable:
they turn a broad coefficient search into exact low-dimensional sign tests.
For weights of higher odd degree the OU polynomial closure remains finite, but
new angular sectors and larger positivity problems appear.  Whether the
matched-equilibrium ordering persists there is left open.

\section{Computational reproducibility and AI assistance}
\label{sec:reproducibility}

The accompanying package contains the exact symbolic Poisson calculation,
complete rational hold-outs, the independent covariance/Wick implementation,
the benchmark reconstruction, figure-generation code, environment versions,
and reference outputs.  All calculations are deterministic.  The exact
residual tests and positive-domain formulas, rather than numerical precision,
are the basis of the result.

OpenAI Codex using the model identifier \texttt{gpt-5.6-sol} provided
substantive assistance with scientific reasoning, symbolic and numerical code,
literature synthesis, verification design, figure generation, and manuscript
drafting.  It was directed to retain the full parameter domain, matched
stiffness, stationary boundary conditions, units, and factor-of-two convention.
Its mathematical output was checked by exact zero-residual identities,
positive-coefficient certificates, complete rational hold-outs, and a
separately implemented covariance/Wick calculation as described above.  These
machine checks do not replace author judgment or accountability; the named
author must complete a final scientific, textual, rights, and policy review
before any submission.

\section*{Data availability}

No experimental data were used.  All code and generated numerical data needed
to reproduce the formulas, checks, and figure are included with this manuscript
in the reproducibility package described in Sec.~\ref{sec:reproducibility}.
Downloaded third-party articles are not part of that package.

\appendix

\section{Completeness of the gradient quotient}
\label{app:quotient}

For a homogeneous vector field of degree $d$ in two variables, the vector
space has dimension $2(d+1)$.  Gradients of $\cH_{d+1}$ form a subspace of
dimension $d+2$: a homogeneous polynomial of positive degree cannot have
identically zero gradient.  The scalar curls have degree $d-1$ and span a
space of dimension $d$.  Since
\begin{equation}
 2(d+1)-(d+2)=d,
\end{equation}
the curl map induces an isomorphism from vector fields modulo gradients to
$\cH_{d-1}$.  Equation~\eqref{eq:curl} supplies an explicit inverse on the
complement $pRx$.  Applying this argument for $d=1$ and $d=3$ gives the
four representatives in Eq.~\eqref{eq:quotient-representative} and proves that
no extensive position-current direction is lost.

For completeness, if $g=\grad\phi$ and $f$ is any polynomial current rate,
stationarity gives
\begin{equation}
 \Var[\phi(x_t)-\phi(x_0)]=O(1),\qquad \Var J_f(t)=O(t).
\end{equation}
Therefore
\begin{equation}
 \frac{|\Cov(J_f(t),\phi(x_t)-\phi(x_0))|}{t}
 \le \frac{O(\sqrt t)}{t}\longrightarrow0.
\end{equation}
This proves invariance of the entire asymptotic bilinear form, not merely its
diagonal values.

\section{Complete symmetry-reduced polynomial sectors}
\label{app:closure}

Set
\begin{equation}
 \begin{aligned}
 X_0&=|X|^2,& V_0&=|V|^2,\\
 C_0&=X\cdot V,& \cA&=X_1V_2-X_2V_1.
 \end{aligned}
\end{equation}
A linearly independent basis of all nonconstant rotational invariants through
total degree four is
\begin{multline}
 X_0,V_0,C_0,\cA, X_0^2,V_0^2,X_0V_0,X_0C_0,X_0\cA,\\
 V_0C_0,V_0\cA,C_0^2,C_0\cA.
\end{multline}
The identity $\cA^2=X_0V_0-C_0^2$ removes the only redundant quartic.  Both
generators in Eq.~\eqref{eq:generators} preserve this 13-dimensional sector.
For the targets $(\cA,X_0\cA)$, the driven stationary means are
\begin{equation}
 \langle\cA\rangle_b=\frac{2b}{a},\qquad
 \langle X_0\cA\rangle_b=\frac{8b}{a^2},
\end{equation}
and both equilibrium means vanish.

For the angular-two sector introduce independent formal complex variables
$z=X_1+iX_2$, $q=V_1+iV_2$ and their conjugates.  The degree-two basis is
$(z^2,zq,q^2)$.  The eight degree-four monomials are
\begin{equation}
 z^j q^{3-j}\bar z,\quad z^j q^{3-j}\bar q,\qquad j=0,1,2,3.
\end{equation}
They are precisely the degree-two and degree-four monomials of rotational
weight $+2$.  The target $z^2\cA$ lies in their span after application of
the generator, and its stationary mean is zero.  Direct substitution of each
solved polynomial $h$ in Eq.~\eqref{eq:poisson} gives an identically zero
residual.  Gaussian moments are evaluated recursively by pairings from the
covariances in Eq.~\eqref{eq:cov-blocks}.

\onecolumngrid
\section{Positive polynomials}
\label{app:polynomials}

The isotropic determinant in Eq.~\eqref{eq:isotropic-components} contains
\begin{align}
P_0(a,u)={}&448a^5+2944a^4u+1700a^4+5568a^3u^2+8220a^3u+2683a^3\notag\\
&+4096a^2u^3+12480a^2u^2+8432a^2u+1872a^2\notag\\
&+1024au^4+6144au^3+8256au^2+3424au+432a\notag\\
&+768u^4+2304u^3+1440u^2+336u+27.
\label{eq:P0}
\end{align}
Its 20 integer coefficients are positive.  The angular coefficient in
Eq.~\eqref{eq:angular-coefficient} contains
\begin{align}
P_2(a,u)={}&10496a^7+60928a^6u+41024a^6
+147200a^5u^2+235488a^5u+70512a^5\notag\\
&+189440a^4u^3+529920a^4u^2+359932a^4u+65772a^4\notag\\
&+136960a^3u^4+591680a^3u^3+702445a^3u^2
+278523a^3u+28404a^3\notag\\
&+52736a^2u^5+335040a^2u^4+633102a^2u^3+458033a^2u^2
+105507a^2u+3564a^2\notag\\
&+8448au^6+83424au^5+251037au^4+299256au^3
+142527au^2+19332au\notag\\
&+4608u^6+30960u^5+61749u^4+50670u^3+16893u^2+1620u.
\label{eq:P2}
\end{align}
All 33 coefficients are positive.  Thus $P_0(a,u)>0$ and $P_2(a,u)>0$
for the entire open domain $a,u>0$, with no numerical lower-bound argument.

\twocolumngrid
\bibliography{manuscript/references}

\end{document}
